\documentclass[11pt]{article}

\usepackage[margin=1in]{geometry}
\usepackage{amsmath,amssymb,amsthm,mathtools}
\usepackage[T1]{fontenc}
\IfFileExists{lmodern.sty}{\usepackage{lmodern}\usepackage{microtype}}{\usepackage[expansion=false]{microtype}}
\usepackage{booktabs}
\usepackage[hidelinks]{hyperref}

\newtheorem{theorem}{Theorem}
\newtheorem{lemma}[theorem]{Lemma}
\newtheorem{proposition}[theorem]{Proposition}
\newtheorem{conjecture}[theorem]{Conjecture}
\newtheorem{remark}[theorem]{Remark}
\newtheorem{definition}[theorem]{Definition}
\newtheorem{corollary}[theorem]{Corollary}

\newcommand{\R}{\mathbb{R}}
\newcommand{\C}{\mathbb{C}}
\newcommand{\Z}{\mathbb{Z}}
\newcommand{\T}{\mathbb{T}}
\newcommand{\eps}{\varepsilon}
\newcommand{\Hup}{\mathbb{H}}
\newcommand{\Sk}{\mathcal S}

\title{The Base-Derivative of the Kneser Family:\\
Selection Principle, Uniqueness, and Existence\\
via a Riemann--Hilbert Problem\\[2mm]
\large Paper V of the mixed-base tetration series}
\author{Janis Justus}
\date{July 2026}

\begin{document}
\maketitle

\begin{abstract}
Let $A_b=\operatorname{slog}_b$ be the regular (Kneser/Paulsen--Cowgill)
superlogarithm family, $\beta=\ln\ln b$, and let
$V=\partial_\beta A_b|_{b=d}$ be its base derivative (the response kernel
of the mixed-base phase law).  We record the selection principle that
characterizes $V$: the cohomological equation
$V(d^y)-V(y)=-A_d'(d^y)\,y\,(\ln d)\,d^y$, reality on $\R$, the prescribed
fixed-point singularity with pole coefficient $dL/d\beta=\ln d\,L^2/(1-\lambda)$
and logarithmic coefficient $d\lambda/d\beta$, and the normalization
$V(1)=0$.  We prove \emph{uniqueness} under these conditions, reduce
\emph{existence} to a scalar Riemann--Hilbert (conjugate-function) problem
for the $\beta$-derivative of the Kneser theta map, and then \emph{solve}
that boundary problem.  The solution rests on three elementary structural
facts established here: the geometric field $G:=\partial_\beta\chi/\chi'$
built from the Koenigs data satisfies the \emph{exact cocycle}
$G(d^y)=(\ln d)\,d^y\,(G(y)-y)$, so that $A_d'\,G+h\circ A_d$ solves the
cohomological equation for \emph{every} $1$-periodic $h$; the boundary
data $u=\operatorname{Im}\bigl[(A_d'\,G)\circ T_d\bigr]$ is $1$-periodic,
extends \emph{continuously} to the closed circle with $u(\Z)=0$ and Dini
modulus $O(\delta\ln^2(1/\delta))$ --- the $\ln\ln$ cascade of the Koenigs
coordinate at the punctures is annihilated by the collapse
$|1/\chi'|=O(\delta\ln(1/\delta))$; and the one-sided limits of the
geometric field at the punctures are the exact orbit values
$V(T_d(m))$, e.g.\ $-e/T_e'(1)$ at $m=1$ for $d=e$ (verified to ten
digits).  Consequently the conjugate-function step is classical
(Privalov--Dini; no weighted class is needed), and interior recovery
follows by Schwarz reflection, Riemann removability at the punctures, and
the $\Sk_\beta/\Sk$ pole expansion at the fixed point.  The outcome is an
existence-and-uniqueness theorem for the selection principle under the
classical branch-geometry package of the Kneser solution, with the
identification $V=\partial_\beta A_b$ under Paulsen's base holomorphy.
The construction is exactly the graded-mesh solver already implemented,
which reproduces the independently measured kernel to $17.8$ digits in
the first Fourier mode and $20.0$ digits in the mean.
\end{abstract}

\section{Setting and the selection principle}

Fix $d>e^{1/e}$, let $T_d$ be the regular tetration branch
($T_d(0)=1$, $T_d(x+1)=d^{T_d(x)}$, increasing from $(-2,\infty)$ onto
$\R$), $A_d=T_d^{-1}$, and let $L$ denote the upper primary fixed point of
$z\mapsto d^z$, with multiplier $\lambda=L\ln d$, $|\lambda|>1$.  For the
base family write $\beta=\ln\ln b$; by Paulsen's base holomorphy
\cite{p19} the family $b\mapsto A_b(y)$ is differentiable in $\beta$ for
$y$ in the real domain, and we set
\[
  V(y)\ :=\ \partial_\beta A_b(y)\big|_{b=d}.
\]
Differentiating the Abel equation $A_b(b^y)=A_b(y)+1$ in $\beta$ gives the
\emph{cohomological equation}
\begin{equation}\label{eq:coh}
  V(d^{\,y})-V(y)\ =\ g_{\mathrm{inh}}(y)\ :=\ -A_d'(d^{\,y})\,y\,(\ln d)\,d^{\,y}.
\end{equation}
Differentiating the Kneser expansion
$A_d(y)=\ln(y-L)/\ln\lambda+\mathcal O(1)$ at $L$ in $\beta$, with
\begin{equation}\label{eq:dL}
  \frac{dL}{d\beta}=\frac{\ln d\;L^2}{1-\lambda},
  \qquad
  \frac{d\lambda}{d\beta}=\ln d\Bigl(\frac{dL}{d\beta}+L\Bigr),
\end{equation}
yields the \emph{prescribed singular part}
\begin{equation}\label{eq:S}
  S(y)\ :=\ -\frac{dL/d\beta}{\ln\lambda\,(y-L)}
          \ -\ \frac{d\lambda/d\beta}{\lambda(\ln\lambda)^2}\,\ln(y-L),
\end{equation}
together with its conjugate at $\bar L$.  Finally $A_b(1)=0$ for every $b$
forces $V(1)=0$.

\begin{definition}[Selection principle]\label{def:selection}
$V$ is the unique function, holomorphic on the domain of $A_d$, satisfying
(i) equation \eqref{eq:coh}; (ii) $V(\R)\subset\R$; (iii) $V-S$ is bounded
near $L$ in the radial principal-sheet sense, with the conjugate structure
at $\bar L$; (iv) $V(1)=0$.
\end{definition}

\subsection*{Standing hypotheses (branch-geometry package)}

The existence proof uses the following classical structure of the regular
branch; we state it explicitly so that the logical status of each step is
transparent.

\begin{itemize}
\item[(K1)] \emph{Strip geometry.}  $T_d$ is holomorphic on the upper
half-plane $\Hup$, continuous up to $\R$, maps $\Hup$ into the attracting
basin of the backward dynamics avoiding the singular backward orbit of
$0$ (so that the principal-log backward orbit, and with it the Koenigs
data of Section~\ref{sec:data}, is defined and holomorphic on
$T_d(\Hup)$), and $T_d(z)\to L$ as $\operatorname{Im}z\to+\infty$.
$A_d$ is holomorphic on a neighbourhood of $T_d\bigl((-2,\infty)\bigr)$
in $\C$ and on $T_d(\Hup)$.
\item[(K2)] \emph{Theta structure.}  With $\chi=\ln\Sk/\ln\lambda$ the
Koenigs--Abel coordinate (Section~\ref{sec:data}), $p:=A_d\circ
T_{\mathrm{reg}}-\mathrm{id}$ is $1$-periodic and holomorphic on
$\{\operatorname{Im}\chi>h_1\}$ for some $h_1$, with $p'\to0$ as
$\operatorname{Im}\chi\to+\infty$ (Kneser's theta decay).
\item[(K3)] \emph{Base holomorphy} (only for the identification):
$b\mapsto A_b(y)$ is holomorphic near $b=d$, locally uniformly in $y$
\cite{p19}.
\end{itemize}

(K1)--(K2) encode the geometry of Kneser's construction \cite{kneser};
they involve the single base $d$ only.  (K3) is used solely to identify
the solution produced below with $\partial_\beta A_b$; existence and
uniqueness do not require it.

\section{Uniqueness}

\begin{theorem}[Uniqueness]\label{thm:unique}
At most one function satisfies (i)--(iv) of
Definition~\ref{def:selection}.
\end{theorem}

\begin{proof}
Let $V_1,V_2$ be two solutions and $h_0:=V_1-V_2$.  By (i), $h_0$ solves
the homogeneous equation $h_0(d^y)=h_0(y)$, hence $h_0=h\circ A_d$ for a
$1$-periodic $h$ on the range of $A_d$, holomorphic where $A_d$ is (define
$h:=h_0\circ T_d$ on a period and extend by periodicity; the functional
equation makes the extension consistent).  Near $L$, the Kneser
superlogarithm satisfies $\operatorname{Im}A_d\to+\infty$ radially on the
principal sheet, and near $\bar L$, $\operatorname{Im}A_d\to-\infty$.
Expand $h(z)=\sum_{k\in\Z}h_k e^{2\pi ikz}$.  By (iii) both $V_i$ have the
same singular part at $L$; hence $h_0=h\circ A_d$ is bounded near $L$,
which forces $h_k=0$ for all $k<0$ (each such mode grows like
$e^{2\pi|k|\operatorname{Im}A_d}$).  The conjugate condition at $\bar L$
forces $h_k=0$ for all $k>0$.  Thus $h$ is a constant, real by (ii), and
$V_1(1)=V_2(1)=0$ gives $h=0$.
\end{proof}

\begin{remark}[Radial versus orbital reading of (iii)]\label{rem:radial}
Condition (iii) must be read radially on the principal sheet.  Along the
backward principal-logarithm spiral $w_{k+1}=\log_d w_k\to L$, the
analytic continuation of $V$ has pole coefficient
$-\bigl(1+p'(\chi)\bigr)\,\frac{dL/d\beta}{\ln\lambda}$, where
$\chi=\ln\mathcal S/\ln\lambda$ is the Koenigs--Abel coordinate and
$p:=A_d-\chi$ the theta correction: the theta derivative multiplies the
pole ($|p'|\approx0.5$ for $d=e$).  An anchor built with the bare $S$ along
the spiral therefore diverges like $|\lambda|^{n}$; this was verified
numerically and is the reason the selection principle cannot be
implemented as a naive orbit summation.
\end{remark}

\section{Two structural lemmas}

\begin{lemma}[Gauge triviality of the natural particular solution]
\label{lem:trivial}
In the Abel coordinate write $U(z):=\partial_\beta T_b(z)|_{b=d}$, which
satisfies $U(z+1)=\ln d\;T_d(z+1)\,\bigl(U(z)+T_d(z)\bigr)$.  The unique
solution decaying up the tower,
\[
  U_{\mathrm{ser}}(z)\ =\ -\sum_{j\ge0}
  \frac{T_d(z+j)}{\prod_{i=1}^{j}\bigl(\ln d\;T_d(z+i)\bigr)},
\]
is hyper-convergent, and the associated kernel candidate vanishes
identically:
\[
  V_{\mathrm{ser}}(T_d(z))\ :=\ -\,\frac{U_{\mathrm{ser}}(z)}{T_d'(z)}
  \ =\ \sum_{j\ge0}\frac{T_d(z+j)}{T_d'(z+j)},
  \qquad\text{hence}\qquad
  G_{\mathrm{ser}}\ \equiv\ 0
\]
for the induced periodic kernel $G(\theta)=V(T_d(\theta))-\sum_{j\ge0}
T_d(j{+}\theta)/T_d'(j{+}\theta)$.
\end{lemma}

\begin{proof}
Chain rule: $T_d'(z+j)=T_d'(z)\prod_{i=1}^{j}\ln d\,T_d(z+i)$, so the
$j$-th summand of $-U_{\mathrm{ser}}/T_d'$ is exactly
$T_d(z+j)/T_d'(z+j)$; the subtraction defining $G$ removes the entire sum.
\end{proof}

Lemma~\ref{lem:trivial} is the precise form of the warning that ``solving
the equation teaches nothing'': every Fourier coefficient of the physical
kernel lies in the homogeneous component that the selection principle must
pin down.

\begin{lemma}[Reduction to a boundary problem]\label{lem:reduction}
Set $g(z):=V(T_d(z))-V_{\mathrm{geo}}(z)$, where
$V_{\mathrm{geo}}:=(1+p'(\chi))\,\partial_\beta\chi$ is computable from the
Koenigs data alone via the backward-orbit recursions
\[
  v_{k+1}=\frac{v_k}{w_k\ln d}-w_{k+1}\ \ (v_0=0,\ v_k\to dL/d\beta),\qquad
  s_{k+1}=\frac{s_k}{w_k\ln d},
\]
\[
  \ln\mathcal S=\lim_n\,[\,n\ln\lambda+\log(w_n-L)\,],\qquad
  \partial_\beta\ln\mathcal S=\lim_n\,[\,n\,d\ln\lambda+(v_n-dL/d\beta)/(w_n-L)\,],
\]
and $1+p'(\chi)=\ln\lambda\,/\,(T_d'\cdot\mathcal S'/\mathcal S)$.
Then $g=p_\beta\circ\chi\circ T_d$ with $p_\beta=\partial_\beta p$ the
base derivative of the theta map: $g$ is $1$-periodic, holomorphic and
bounded in the upper half-plane with a limit at $i\infty$, and continuous
on $\R\setminus\Z$.  The Koenigs coordinate itself has the explicit
$\ln\ln$-endpoint cascade inherited from the singularity of
$\chi$ at the normalization orbit $1\mapsto0\mapsto-\infty$
($\chi(\theta)\sim[\ln(\ln(\tau\theta)-L)]/\ln\lambda$ as
$\theta\downarrow0$, $\tau=T_d'(-1)$; one cascade level deeper at
$\theta\uparrow1$); the boundary \emph{data}, however, is continuous up to
the punctures (Lemma~\ref{lem:data} below).  Conditions (ii)--(iv) become
the scalar Riemann--Hilbert problem
\begin{equation}\label{eq:RH}
  \operatorname{Im}g(\theta)\ =\ -\operatorname{Im}V_{\mathrm{geo}}(\theta)
  \quad(\theta\in\R\setminus\Z),
  \qquad
  g\in H^{+}_{\mathrm{per}},
  \qquad
  g(0)=0 ,
\end{equation}
where $H^{+}_{\mathrm{per}}$ denotes $1$-periodic boundary values of
functions holomorphic and bounded in the upper half-plane, and the
normalization $g(0)=0$ expresses $V(1)=0$ through
$V_{\mathrm{geo}}(0^+)\to0$ (Lemma~\ref{lem:data}).
\end{lemma}

\begin{proof}[Proof sketch]
$A_d=\chi+p(\chi)$ with $p$ $1$-periodic and decaying as
$\operatorname{Im}\chi\to+\infty$ (Kneser); differentiate in $\beta$ and
use $1+p'=dA_d/d\chi=A_d'/\chi'$.  All stated limits exist by the Koenigs
theory; the identities were verified against the $\beta$-Koenigs equation
to $10^{-121}$.  Composition with $\chi\circ T_d$, which maps the upper
half-plane into the domain of $p$ up to the $\Z$-punctures (K1)--(K2),
gives the mapping properties of $g$; boundedness up to the punctures
follows from Lemma~\ref{lem:data} and $g=V\circ T_d-V_{\mathrm{geo}}$
under (K3), and unconditionally from the construction of
Section~\ref{sec:interior}, which re-derives $g$ without reference to
$V$.
\end{proof}

\section{Regularity of the boundary data}\label{sec:data}

This section contains the new structural input that makes the existence
program close: the boundary data of \eqref{eq:RH} is far better behaved
than the $\ln\ln$ cascade of the coordinate $\chi$ suggests.  Everything
rests on an exact cocycle for the \emph{geometric field}
\[
  G\ :=\ \frac{\partial_\beta\chi}{\chi'},
  \qquad\text{so that}\qquad
  V_{\mathrm{geo}}\ =\ (1+p'(\chi))\,\partial_\beta\chi
  \ =\ A_d'\,G
\]
(using $A_d'=(1+p'(\chi))\chi'$).  Here $\chi=\ln\Sk/\ln\lambda$ and
$\partial_\beta\chi$ are the Koenigs--Abel coordinate of the repelling
fixed point $L_b$ and its base derivative; both are elementary objects:

\begin{lemma}[Koenigs base regularity]\label{lem:koenigs-beta}
For $b$ in a complex neighbourhood of $d$, the fixed point
$L_b=-W(-\ln b)/\ln b$ (principal upper branch), the multiplier
$\lambda_b=L_b\ln b$, and the Koenigs linearizer $\Sk_b$ (normalized
$\Sk_b(L_b)=0$, $\Sk_b'(L_b)=1$) depend holomorphically on $b$, locally
uniformly on the basin.  Consequently $\chi_b$, $\partial_\beta\chi$, and
the recursions of Lemma~\ref{lem:reduction} are well defined without any
family input beyond the single-base dynamics.
\end{lemma}

\begin{proof}[Proof sketch]
$L_b$ is holomorphic by the implicit function theorem (or the Lambert-$W$
formula) since $\lambda_b\ne1$.  The Koenigs limit
$\Sk_b(y)=\lim_n\lambda_b^{\,n}\,(\ell_b^{\circ n}(y)-L_b)$ along the
principal-log backward orbit $\ell_b=\log_b$ converges locally uniformly
in $(y,b)$ by the uniform spiral contraction near $L_b$
($|\lambda_b|>1$ locally uniformly), hence is jointly holomorphic
(Weierstrass).  The recursions are the term-wise $\beta$-derivatives of
this limit.
\end{proof}

\begin{lemma}[Exact cocycle for the geometric field]\label{lem:cocycle}
On $T_d(\Hup)$ and on the real orbit domain,
\begin{equation}\label{eq:Gcoc}
  G(d^{\,y})\ =\ (\ln d)\,d^{\,y}\,\bigl(G(y)-y\bigr).
\end{equation}
Consequently, for \emph{every} $1$-periodic $h$ the field
$\widetilde V:=A_d'\,G+h\circ A_d$ satisfies the cohomological equation
\eqref{eq:coh}; and $u:=\operatorname{Im}V_{\mathrm{geo}}
=\operatorname{Im}\bigl[(A_d'G)\circ T_d\bigr]$ is $1$-periodic on
$\R\setminus\Z$.
\end{lemma}

\begin{proof}
From $\chi_b(y)=\chi_b(\log_b y)+1$: differentiate in $\beta$ at fixed
$y$, using $\partial_\beta \log_b y=-\log_b y$ (since
$\partial_\beta\ln b=\ln b$), to get
$\partial_\beta\chi(y)=\partial_\beta\chi(w)-w\,\chi'(w)$ with
$w=\log_d y$; and $\chi'(y)=\chi'(w)/(y\ln d)$ by the chain rule.
Divide.  This is \eqref{eq:Gcoc} written at $y\mapsto d^{\,y}$.
For the second claim, the Abel equation gives
$A_d'(d^{\,y})\,(\ln d)\,d^{\,y}=A_d'(y)$, hence
\[
  A_d'(d^y)G(d^y)-A_d'(y)G(y)
  = A_d'(y)\bigl(G(y)-y\bigr)-A_d'(y)G(y)=-\,y\,A_d'(y)
  = -A_d'(d^{\,y})\,y\,(\ln d)\,d^{\,y},
\]
which is \eqref{eq:coh}; the $h\circ A_d$ part is $1$-periodic in the
Abel coordinate and cancels.  Finally, applying this identity on
$\R\setminus\Z$ and noting that the right-hand side of \eqref{eq:coh} is
\emph{real} there shows that
$V_{\mathrm{geo}}(\theta+1)-V_{\mathrm{geo}}(\theta)\in\R$, i.e.\ $u$ is
$1$-periodic.
\end{proof}

\begin{lemma}[Endpoint regularity of the data]\label{lem:data}
Let $y_m:=T_d(m)$, $m\in\Z_{\ge0}$, be the orbit points of $1$.  Define
$G_0:=0$ and $G_{m+1}:=(\ln d)\,y_{m+1}(G_m-y_m)$, and set
$V_m:=A_d'(y_m)\,G_m$ (so $V_0=0$,
$V_1=-A_d'(d)\,d\ln d=-\,d\ln d/T_d'(1)$).  Then, in every sector
$\{|z-m|<\delta_0,\ \operatorname{Im}z\ge0\}$,
\begin{equation}\label{eq:collapse}
  V_{\mathrm{geo}}(z)\ =\ V_m\ +\
  O\!\bigl(|z-m|\,\ln^2(1/|z-m|)\bigr),
\end{equation}
from both sides and locally uniformly in the closed sector.  In
particular $u=\operatorname{Im}V_{\mathrm{geo}}$ extends to a
$1$-periodic function in $C(\T)$ with $u(\Z)=0$ and Dini modulus
\[
  \omega_u(\delta)\ \le\ C\,\delta\,\bigl(1+\ln(1/\delta)\bigr)^{2}.
\]
\end{lemma}

\begin{proof}
Write $\delta:=|z-m|$ and iterate the cocycle \eqref{eq:Gcoc} backwards:
$G(y)=y^{-1}(\ln d)^{-1}\,G(d^y)+y$ expresses $G$ at a point in terms of
$G$ one backward step down; equivalently, going down the backward orbit
$w_1=\log_d y,\ w_2=\log_d w_1,\dots$ of $y=T_d(z)$,
\[
  G(y)\ =\ (\ln d)\,y\,\bigl(G(w_1)-w_1\bigr),\qquad
  G(w_1)\ =\ (\ln d)\,w_1\bigl(G(w_2)-w_2\bigr),\ \dots
\]
For $z$ near $m$ the first $m{+}1$ orbit points are within $O(\delta)$ of
the finite orbit $y_m,y_{m-1},\dots,y_0=1$, after which the orbit passes
the two singular-entry steps: $w_{m+1}=\log_d(1+O(\delta))=O(\delta)$ and
$w_{m+2}=\log_d w_{m+1}=-t/\ln d+O(1)$ with $t:=\ln(1/\delta)$, and then
$w_{m+3}=\log_d w_{m+2}$, of modulus $O(\ln t)$, lies in a fixed compact
set $K$ of the basin on which finitely many further principal-log steps
reach the spiral annulus of $L$.  On $K$ and on the annulus, $G$ is
holomorphic and bounded (Lemma~\ref{lem:koenigs-beta}: $\chi'\ne0$ there
because $\Sk'\ne0$ along the non-singular backward orbit, and
$\partial_\beta\chi$ is holomorphic), say $|G|\le B$.  Unwinding the
three explicit steps:
\[
  |G(w_{m+2})|\le C\,\ln t\,(B+\ln t)=O(\ln^2 t),\qquad
  |G(w_{m+1})|\le C\,\delta\,\bigl(O(\ln^2 t)+t\bigr)=O(\delta\,t),
\]
\[
  G(T_d(z))\ =\ (\ln d)^{m+1}\,y_m\cdots y_0\,\cdot
  \bigl(G(w_{m+1})-w_{m+1}\bigr)\ \prod\text{-corrections}
  \ =\ G_m+O(\delta\,t)\cdot C_m+O(\delta),
\]
where the finite products over the $m{+}1$ regular steps converge to the
recursion defining $G_m$ with Lipschitz error $O(\delta)$ (all maps
involved are holomorphic with bounded derivative near the finite orbit).
Multiplying by $A_d'(T_d(z))=A_d'(y_m)+O(\delta)$ gives
\eqref{eq:collapse} with the stated $\delta\ln$-type error (one extra
factor $\ln t$ is kept as margin for the $m$-dependence of the entry
step; for $m=0$ the term $G_0-y_0\cdot$-correction is absent and the same
bound holds directly).  The two-sided statement follows because the
argument only used $|w_{m+1}|=O(\delta)$, which holds on both sides
(on the left side of $m$ the small orbit point is negative and the next
step picks up the $+i\pi$ branch, changing only the constants).  Since
$V_m\in\R$ (all $y_m,G_m$ real), the imaginary part of
\eqref{eq:collapse} gives $|u|\le C\delta\ln^2(1/\delta)$ near the
punctures, with $u(m^\pm)\to0$.  The modulus of continuity follows by
combining this with the interior gradient bound
$|u'(\theta)|\le C\ln^2(1/\mathrm{dist}(\theta,\Z))$, obtained from
\eqref{eq:collapse} by the Cauchy estimate on the disc of radius
$\mathrm{dist}(\theta,\Z)/2$ contained in the sector.
\end{proof}

\begin{remark}[Why the $\ln\ln$ cascade is harmless]\label{rem:collapse}
The coordinate $\chi\circ T_d$ does diverge at the punctures through the
iterated-logarithm cascade, and $|\partial_\beta\chi|$ even grows like
$\delta^{-1}$ on the deeper side; but the data of \eqref{eq:RH} is the
\emph{product} $A_d'\,G$, and the cocycle shows the singular factors
cancel to leave the finite orbit values $V_m$.  The original solvability
conjecture (v1 of this note) posited a weighted class accommodating
$\ln\ln$ tails; the correct statement is stronger: the data is continuous
with a Dini modulus, and \emph{no weighted class is needed}.  It is also
worth noting how slowly the cascade actually moves: the depth of the
iterated-log cascade increases by one only when $\ln(1/\delta)$ is
exponentiated, so $\chi\circ T_d$ stays in a compact region for all
$\delta>10^{-10^{17}}$; but the proof above does not use this, since the
collapse \eqref{eq:collapse} holds uniformly in the depth.
\end{remark}

\begin{proposition}[Numerical confirmation]\label{prop:numcheck}
For $d=e$ (all entries from the exact Koenigs recursions at $140$
digits): the collapse \eqref{eq:collapse} at $m=0$ gives
$|V_{\mathrm{geo}}(\delta)|/(\delta\ln(1/\delta))
=6.2,\,7.7,\,8.8,\,9.7,\,10.4,\,11.1$ at
$\delta=10^{-2},10^{-4},\dots,10^{-12}$ (slow $\ln$-growth, consistent
with the $\delta\ln^2$ bound); at $m=1$,
$V_{\mathrm{geo}}(1-\delta)\to V_1=-e/T_e'(1)=-0.91594605649953\ldots$
with $|V_{\mathrm{geo}}(1-10^{-12})-V_1|\approx3\cdot10^{-10}$, and the
same limit from the right side; the imaginary part decays like
$\delta\ln(1/\delta)$ on both sides
($\operatorname{Im}V_{\mathrm{geo}}(1-\delta)=2.7\cdot10^{-1},\dots,
3.0\cdot10^{-10}$ for $\delta=10^{-2},\dots,10^{-12}$); and the
periodicity of $u$ (Lemma~\ref{lem:cocycle}) holds to print precision,
e.g.\ $u(-10^{-4})=u(1-10^{-4})=0.006667$.
\end{proposition}

\section{Solution of the boundary problem (step 2)}\label{sec:step2}

\begin{theorem}[Solvability of \eqref{eq:RH}]\label{thm:step2}
Let $u\in C(\T)$ be real with a Dini modulus of continuity (as furnished
by Lemma~\ref{lem:data}) and $u(0)=0$.  Then the boundary problem
\eqref{eq:RH} has exactly one solution $g$ in the class
$H^{+}_{\mathrm{per}}\cap C(\overline\Hup\ \mathrm{mod}\ 1)$, namely
\[
  g\ =\ \widetilde u\ -\ i\,u\ -\ \widetilde u(0),
\]
where $\widetilde u$ is the conjugate function
($\widehat{\widetilde u}(k)=-i\,\mathrm{sgn}(k)\,\widehat u(k)$).
Moreover $g$ is continuous on the \emph{closed} circle --- including the
former punctures --- and $g(i\infty)=\widehat u(0)\,(-i)+\ldots$ exists
(the mean of $g$).
\end{theorem}

\begin{proof}
\emph{Existence.}  Since $u$ is Dini-continuous, the conjugate function
$\widetilde u$ exists at every point and is continuous on $\T$
(Privalov's theorem in its Dini form; see \cite[Ch.~III]{garnett} or
\cite[Ch.~III, Thm.~13.30]{zygmund}); moreover $u+i\widetilde u$ is the
boundary value of the holomorphic function
$2\sum_{k\ge1}\widehat u(k)e^{2\pi ik z}+\widehat u(0)$ on $\Hup$, which
extends continuously to $\overline\Hup$ (Fej\'er means converge uniformly
for continuous boundary data with continuous conjugate).  Then
$g:=-i\,(u+i\widetilde u)-\widetilde u(0)=\widetilde u-iu-\widetilde u(0)$
lies in $H^{+}_{\mathrm{per}}\cap C$, has
$\operatorname{Im}g=-u$ on $\R$, and $g(0)=\widetilde u(0)-i\cdot0-
\widetilde u(0)=0$.

\emph{Uniqueness.}  If $g_1,g_2$ are two solutions, $h:=g_1-g_2\in
H^{+}_{\mathrm{per}}$ is bounded with real boundary values.  Its Fourier
expansion on a horizontal line has only modes $k\ge0$ (boundedness at
$i\infty$ kills $k<0$); reality on $\R$ forces
$\overline{h_k}=h_{-k}$, hence $h_k=0$ for $k\ne0$ and $h_0\in\R$; the
normalization $h(0)=0$ gives $h\equiv0$.
\end{proof}

\begin{remark}[Fourier tails]\label{rem:tails}
The Dini modulus $\omega_u(\delta)\le C\delta\ln^2(1/\delta)$ gives
$|\widehat u(k)|\le C\,\omega_u(1/k)\le C\ln^2 k/k$ --- consistent with
the measured $|k|^{-1.28}$-type decay of the mode contributions in the
numerical solver over the fitted window, and fast enough for the
$L^2$- and uniform convergence used above.
\end{remark}

\section{Interior recovery and the singular normalization (step 3)}
\label{sec:interior}

\begin{theorem}[Existence]\label{thm:exist}
Assume (K1)--(K2).  Define, with $u:=\operatorname{Im}V_{\mathrm{geo}}$
from Section~\ref{sec:data} and $g$ from Theorem~\ref{thm:step2},
\[
  \widetilde V\circ T_d\ :=\ V_{\mathrm{geo}}\ +\ g
  \qquad\text{on }\ \overline\Hup .
\]
Then $\widetilde V$, extended to the domain of $A_d$ by the cohomological
equation, satisfies (i)--(iv) of Definition~\ref{def:selection}.  In
particular, together with Theorem~\ref{thm:unique}, the selection
principle has \emph{exactly one} solution.  Under (K3) this solution is
$\partial_\beta A_b|_{b=d}$.
\end{theorem}

\begin{proof}
\emph{Interior holomorphy and reflection.}
$V_{\mathrm{geo}}=A_d'\,G$ is holomorphic on $\Hup$ by (K1) and
Lemma~\ref{lem:koenigs-beta} (the backward orbit of $T_d(z)$ avoids the
singular set for $z\in\Hup$), and $g$ is holomorphic on $\Hup$ by
Theorem~\ref{thm:step2}.  On $\R\setminus\Z$ the boundary values of the
sum are continuous and \emph{real}:
$\operatorname{Im}(V_{\mathrm{geo}}+g)=u+(-u)=0$.  By Schwarz reflection,
$\widetilde V\circ T_d$ extends holomorphically across every interval of
$\R\setminus\Z$.  At a puncture $m\in\Z$, Lemma~\ref{lem:data} (sector
version) and the continuity of $g$ on the closed half-plane
(Theorem~\ref{thm:step2}) show that the reflected function is bounded
near $m$; by Riemann's removability theorem it extends holomorphically
across $m$ as well, with value $V_m+g(m)$.  Hence
$\widetilde V\circ T_d$ is holomorphic on a full neighbourhood of $\R$,
and $\widetilde V=(\widetilde V\circ T_d)\circ A_d$ is holomorphic on a
neighbourhood of the real orbit domain (recall $A_d'\neq0$ on $\R$); the
cohomological equation transports holomorphy to the rest of the domain of
$A_d$.

\emph{(i).}  By Lemma~\ref{lem:cocycle}, $A_d'G$ satisfies \eqref{eq:coh}
and $g=h\circ A_d$-type terms are periodic in the Abel coordinate:
$g(\theta+1)=g(\theta)$, so $\widetilde V$ satisfies \eqref{eq:coh}.

\emph{(ii).}  Boundary values on $\R\setminus\Z$ are real by
construction; by continuity and the removability step, real on all of the
real orbit domain.

\emph{(iii).}  Radially at $L$, i.e.\ $\operatorname{Im}z\to+\infty$ in
the strip picture (K1): expand the Koenigs data at the fixed point.
With $\Sk_b(y)=(y-L_b)(1+O(y-L_b))$ and $\Sk_b'(L_b)=1$,
\[
  \partial_\beta\ln\Sk\ =\ \frac{\Sk_\beta}{\Sk}
  \ =\ -\,\frac{dL/d\beta}{y-L}+O(1),
  \qquad
  \chi=\frac{\ln(y-L)}{\ln\lambda}+O(1),
\]
hence
\[
  \partial_\beta\chi
  \ =\ \frac{\partial_\beta\ln\Sk-\chi\,d\ln\lambda}{\ln\lambda}
  \ =\ -\,\frac{dL/d\beta}{\ln\lambda\,(y-L)}
  \ -\ \frac{d\ln\lambda}{(\ln\lambda)^2}\,\ln(y-L)\ +\ O(1)
  \ =\ S(y)+O(1),
\]
which is exactly \eqref{eq:S} (note $d\ln\lambda=(d\lambda/d\beta)/
\lambda$).  The theta factor satisfies $p'(\chi)\to0$ as
$\operatorname{Im}\chi\to+\infty$ by (K2), and
$\operatorname{Im}\chi\sim-\operatorname{Im}(1/\ln\lambda)\cdot
\ln\!\frac1{|y-L|}\to+\infty$ radially, with
$|p'(\chi)|=O(|y-L|^{\,2\pi\nu})$, $\nu:=-\operatorname{Im}(1/\ln\lambda)
>0$ (for $d=e$: $2\pi\nu=4.447$); hence
$(1+p'(\chi))\partial_\beta\chi=S(y)+O(1)$ radially.  Finally $g$ is
bounded on $\Hup$ with a limit at $i\infty$
(Theorem~\ref{thm:step2}), so $\widetilde V-S$ is bounded radially at
$L$.  The conjugate structure at $\bar L$ holds by reflection (reality on
$\R$).

\emph{(iv).}  $\widetilde V(1)=\lim_{\theta\downarrow0}
\bigl[V_{\mathrm{geo}}(\theta)+g(\theta)\bigr]=0+g(0)=0$ by
Lemma~\ref{lem:data} and the normalization of Theorem~\ref{thm:step2}.

\emph{Identification.}  Under (K3), $V=\partial_\beta A_b$ exists and
satisfies (i)--(iv) (Section~1); by Theorem~\ref{thm:unique},
$V=\widetilde V$.
\end{proof}

\begin{corollary}[Status of the v1 conjecture]\label{cor:conj}
The Solvability Conjecture of the first version of this note holds, in
the stronger form of Theorems~\ref{thm:step2}--\ref{thm:exist}: the
weighted class is not needed, the solution lies in
$H^{+}_{\mathrm{per}}\cap C$, and the cascade subtraction step (1) of the
proposed route is superseded by the collapse identity
\eqref{eq:collapse}.  Consequently the selection principle characterizes
the response kernel under the classical branch geometry (K1)--(K2) alone,
and base differentiability of the Kneser family on the real domain, with
the explicit kernel, follows for any family satisfying (K3).
\end{corollary}

\begin{remark}[Honest scope]\label{rem:scope}
What remains outside this note is classical: (K1)--(K2) are the geometry
of Kneser's Riemann-mapping construction \cite{kneser} (holomorphy of the
branch in the upper half-plane, approach to $L$ at $i\infty$, theta decay)
--- properties of the single base $d$, independent of the base family;
their verification for the regular branch is the content of the classical
construction and of the companion paper's Appendix on stopped
coordinates, and all quantitative consequences used here were verified
numerically far beyond doubt (Proposition~\ref{prop:numcheck};
$\beta$-Koenigs gates to $10^{-121}$).  The proofs of
Sections~\ref{sec:data}--\ref{sec:interior} are self-contained given that
package.
\end{remark}

\paragraph{Numerical realization (constructive check).}
The problem \eqref{eq:RH} was solved directly on a graded mesh: composite
Gauss--Legendre panels on geometric cascades at both endpoints, data
$u=\operatorname{Im}V_{\mathrm{geo}}$ evaluated pointwise by the exact
Koenigs recursions, periodic Hilbert transform with exact cotangent-kernel
integrals, and normalization through the periodic kernel at $\theta=0$
(the non-periodic $V\circ T_d$ must not be trigonometrically
interpolated).  This is precisely the constructive scheme of
Theorems~\ref{thm:step2}--\ref{thm:exist}.  A convergence study (panel
order $24\to48$, decade $\to$ half-decade grading, truncation
$10^{-30}\to10^{-38}$) shows the expected spectral behaviour: for $d=e$
the engine-free construction reproduces the independently measured kernel
($\eps$-stencils of the base family, ``lnln-exact'' parametrization) to
\[
  \text{$34.1$ digits in }\ \bar G=\smash{\int_0^1}G,\qquad
  \text{$34.4$ digits in }\ |\widehat G(1)|,\qquad
  \text{$34.1$ digits in }\ |\widehat G(2)|,
\]
i.e.\ to the full stored length of the reference, with internal
self-convergence $2.9\cdot10^{-38}$ (relative, in $\bar G$) between the
two finest meshes.  Across the base family $d\in\{2,e,3,5,10\}$ the
solver converges uniformly ($\approx29.7/27.7/25.5$ digits already at
panel order $40$ with decade grading, $\gtrsim30$--$34$ at the sharpened
settings); for the extreme bases $d=2$ and $d=10$ the RH construction
and the $\eps$-stencil route agree to $30$--$34$ digits and jointly
expose an error in the earlier complex-step reference values (correct to
only $\approx20$ resp.\ $\approx15$ digits there --- the engine floor at
extreme bases: $|\lambda(2)|=1.23$ near the parabolic edge, large
$\ln d$ narrowing the theta strip at $d=10$).  The accuracy is set by
the quadrature order and endpoint truncation alone; the formulation
carries at least $34$ digits.

\begin{remark}[What this buys downstream]
With $V$ characterized, the linear-response description of the mixed-base
phase family ($\widehat P_k(b,d)=\eps\widehat G_k(d)+\eps^2\widehat H_k(d)
+O(\eps^3)$, verified family-wide) rests on a selection principle rather
than on measurements; the rank-one inner-base law of the main paper
becomes its leading order; and the base-differentiability building block
requested by Remark~7.4 of the companion diffeomorphism note is in place,
now as a theorem under the classical branch geometry.
\end{remark}


\section*{Authorship and AI Contribution Statement}
This paper is part of a five-paper series produced in a directed human--AI
collaboration. The author conceived the mixed-base phase program, built and
maintains the high-precision tetration engine used for every computation in
the series (a fork of Levenstein's \texttt{fatou.gp}; engine, data and
scripts at \url{https://github.com/Lightrunnerwastaken/gp-tetration}), and directed the research: posing the questions,
setting priorities between rounds, auditing results, and deciding which
claims met the acceptance gates. Substantial parts of the mathematical
development and of the numerical work were produced by AI systems under
that direction: Anthropic's Claude in a theory and review role, an
autonomous research agent operating the numerical pipeline, and OpenAI's
ChatGPT, which contributed to the original versions of Papers I and II and
to manuscript review. In line with prevailing journal policies, AI systems
are not listed as authors; the author assumes full and sole responsibility
for all statements and proofs, and readers should weigh this provenance in
their assessment. The five papers of the series are permanently archived on Zenodo under the concept DOIs \url{https://doi.org/10.5281/zenodo.21346803} (I), \url{https://doi.org/10.5281/zenodo.21361967} (II), \url{https://doi.org/10.5281/zenodo.21362206} (III), \url{https://doi.org/10.5281/zenodo.21363593} (IV), and \url{https://doi.org/10.5281/zenodo.21363809} (V).

\begin{thebibliography}{9}
\bibitem{paper} J.~Justus, \emph{A Phase Law for Mixed-Base Tetration}, Paper~I of the mixed-base tetration series, preprint, July 2026. \url{https://github.com/Lightrunnerwastaken/gp-tetration}. DOI: \url{https://doi.org/10.5281/zenodo.21346803}.
\bibitem{unified} J.~Justus, \emph{The Mixed-Base Tetration Phase Map Is a Circle Diffeomorphism}, Paper~II of the mixed-base tetration series, preprint, July 2026. \url{https://github.com/Lightrunnerwastaken/gp-tetration}. DOI: \url{https://doi.org/10.5281/zenodo.21361967}.
\bibitem{stopp} J.~Justus, \emph{Stopped Channels and Offset Alignment}, Paper~III of the mixed-base tetration series, preprint, July 2026. \url{https://github.com/Lightrunnerwastaken/gp-tetration}. DOI: \url{https://doi.org/10.5281/zenodo.21362206}.
\bibitem{kneser} H.~Kneser, \emph{Reelle analytische L\"osungen der
Gleichung $\varphi(\varphi(x))=e^x$ und verwandter Funktionalgleichungen},
J.~reine angew.~Math.\ 187 (1950), 56--67.
\bibitem{p19} W.~Paulsen, \emph{Tetration for complex bases}, Adv.\
Comput.\ Math.\ 45 (2019).
\bibitem{garnett} J.~Garnett, \emph{Bounded Analytic Functions}, Springer
GTM 236, 2007.
\bibitem{zygmund} A.~Zygmund, \emph{Trigonometric Series}, 3rd ed.,
Cambridge, 2002.
\end{thebibliography}

\end{document}
